\section{Protocolos de transporte de red}
\label{sec:protocolos_red}

Los protocolos de aplicación que gobiernan la contribución y la distribución audiovisual
se construyen sobre un conjunto reducido de protocolos de transporte de red. Comprender
su comportamiento es fundamental para interpretar las decisiones de diseño que subyacen
a cada solución de producción, ya que sus propiedades determinan directamente la latencia,
la fiabilidad y la escalabilidad del sistema.

\subsection{UDP (\textit{User Datagram Protocol})}

UDP es un protocolo de transporte sin conexión definido en la RFC~768~\cite{rfc768_udp}.
Envía datagramas de forma independiente sin establecer sesión previa, sin confirmación de
entrega, sin reordenación de paquetes y sin control de congestión. Estas características
le confieren la mínima latencia posible dentro de la pila de red, al eliminar todo el
\textit{overhead} de fiabilidad.

En producción audiovisual, UDP es el protocolo de elección cuando la latencia prima sobre
la integridad de los datos. Un fotograma perdido es preferible a un fotograma retrasado en
aplicaciones de tiempo real. Por este motivo, los siguientes protocolos de nivel superior
operan sobre UDP:

\begin{itemize}
  \item \textbf{RTP} (\textit{Real-time Transport Protocol}): encapsula los flujos
        multimedia con marcas de tiempo y números de secuencia que permiten al receptor
        detectar y compensar pérdidas sin retransmisión.
  \item \textbf{WebRTC}: emplea UDP con DTLS/SRTP para comunicaciones en tiempo real
        con latencias de 100--500~ms, asumiendo pérdidas ocasionales de paquetes.
  \item \textbf{SRT}: utiliza UDP como transporte pero añade una capa ARQ
        (\textit{Automatic Repeat Request}) propia para recuperar paquetes perdidos
        sin incrementar la latencia en condiciones normales de red.
\end{itemize}

\subsection{TCP (\textit{Transmission Control Protocol})}

TCP es un protocolo de transporte orientado a conexión definido en la
RFC~793~\cite{rfc793_tcp}. Garantiza la entrega ordenada de todos los bytes mediante
confirmaciones de recepción, retransmisión automática de segmentos perdidos y control
de flujo y congestión. Estas garantías introducen latencia variable en redes con pérdidas,
ya que el receptor debe esperar la retransmisión antes de procesar datos posteriores.

En producción audiovisual, TCP se emplea cuando la integridad de los datos no es
negociable y la latencia adicional es aceptable:

\begin{itemize}
  \item \textbf{RTMP}: protocolo de ingesta basado en TCP que prioriza la entrega
        completa de cada fotograma sobre plataformas como YouTube Live o Twitch.
  \item \textbf{Voctomix interno}: el núcleo voctocore acepta las fuentes de cámara
        mediante conexiones TCP que transportan flujos MKV con vídeo RAW~I420 y audio
        PCM~S16LE. TCP garantiza la sincronización audio-vídeo sin pérdidas dentro
        de la red local.
  \item \textbf{Protocolo de control}: el servidor de comandos de voctocore escucha
        en el puerto 9999 mediante TCP, asegurando que cada comando enviado por la
        interfaz gráfica llega íntegramente al núcleo y que la respuesta de estado
        retorna sin corrupción.
\end{itemize}

\subsection{HTTP / HTTPS (\textit{Hypertext Transfer Protocol})}

HTTP es un protocolo de capa de aplicación definido en la RFC~9110~\cite{rfc9110_http}
que opera sobre TCP (HTTP/1.1 y HTTP/2) o sobre QUIC/UDP (HTTP/3). Su modelo de
petición-respuesta y su infraestructura de caché y distribución global mediante CDN
(\textit{Content Delivery Networks}) lo convierten en el protocolo de referencia para
la distribución de contenidos a audiencias masivas.

En producción audiovisual, HTTP actúa como capa de transporte para los principales
protocolos de distribución adaptativa:

\begin{itemize}
  \item \textbf{HLS y MPEG-DASH}: dividen el vídeo en segmentos de fichero servidos
        como respuestas HTTP estándar. Cualquier infraestructura de servidores web
        o CDN puede distribuir estos segmentos sin configuración específica de media.
  \item \textbf{API REST de telemetría}: el servicio de telemetría de Voctomix~2.0
        expone el estado de la sesión mediante HTTP en el puerto 8080, y el módulo
        de exportación de la GUI envía actualizaciones de estado mediante
        peticiones HTTP POST.
\end{itemize}

\begin{table}[H]
  \centering
  \caption{Comparativa de los protocolos de transporte en producción audiovisual.}
  \label{tab:protocolos_red}
  \small
  \begin{tabularx}{\linewidth}{|l|c|c|c|X|}
    \hline
    \textbf{Protocolo} & \textbf{Orientado a conexión} & \textbf{Fiabilidad} & \textbf{Latencia} & \textbf{Uso en Voctomix} \\
    \hline
    UDP  & No  & Sin garantías & Mínima       & RTP, SRT, WebRTC \\
    TCP  & Sí  & Garantizada   & Variable     & MKV/fuentes, RTMP, control \\
    HTTP & Sí* & Garantizada   & Alta         & HLS, DASH, telemetría REST \\
    \hline
  \end{tabularx}
  \begin{flushleft}
    \footnotesize *HTTP/1.1 y HTTP/2 operan sobre TCP; HTTP/3 opera sobre QUIC (UDP).
  \end{flushleft}
\end{table}
